\documentclass[10pt,a4paper]{amsart}
\usepackage[margin=19mm]{geometry}
\usepackage{amsmath,amssymb,mathtools}
\usepackage[scr=rsfs,cal=euler]{mathalfa}
\usepackage{xfrac}
\usepackage[unicode,hidelinks,bookmarks=false]{hyperref}
\newcommand{\ie}{i.e.\ }
\newcommand{\cf}{cf.\ }
\newcommand{\resp}{respectively\ }
\newcommand{\Subsection}{\subsection}
\providecommand{\nobreakdash}{\nobreak\hbox{-}}
\setlength{\emergencystretch}{2em}
\multlinegap0pt
\usepackage{bm}
\usepackage{bbm}
\usepackage{dsfont}
\usepackage{xspace}
\usepackage{cancel}
\usepackage{mathtools}
\usepackage{amsthm}
\makeatletter
\@ifundefined{theorem}{\newtheorem{theorem}{Theorem}}{}
\@ifundefined{lemma}{\newtheorem{lemma}[theorem]{Lemma}}{}
\makeatother
\title[Directional $\rho$-coefficients]{Directional $\rho$-coefficients}
\author{Enrique de Amo, David Garc\'ia-Fern\'andez, Manuel \'Ubeda-Flores}
\date{Source: arXiv:2505.22206v1 (CC BY 4.0)}
\begin{document}
\maketitle

\noindent\emph{Source and licence.} Directional $\rho$-coefficients. Version arXiv:2505.22206v1 (CC BY 4.0), distributed under the Creative Commons Attribution 4.0 International licence, as recorded in the arXiv licence metadata for this exact version and shown on the abstract page https://arxiv.org/abs/2505.22206v1. Full rights notice: licenses/arxiv-2505.22206-CC-BY-4.0.txt.\par\smallskip

\noindent\textit{Editorial scope.} Counted unit: the Theorem whose source label is \texttt{th:est} in the article (source0.tex lines 582--597, source label \texttt{th:est}), delivered together with the article's own complete original proof of it (source0.tex lines 599--640, one \texttt{proof} environment of 42 lines). No mathematics is written, repaired or re-derived: statement and proof are the article's own text line for line. Required context retained uncounted: lem:aux (lines 556--562). Every definition, display and label of those lines is retained unchanged. Omitted from this package and not redistributed: 1--555, 563--581, 598--598, 641--1031. Nothing inside a retained excerpt is omitted or altered: every retained body is delivered exactly as the article's lines are written, so no omission receipt of any kind accompanies this package. Citations: the retained excerpts carry no citation command at all, so no bibliography entry of the article is transcribed and no reference is redistributed. Reference adaptation: none was needed and none was made; every reference inside the delivered unit resolves inside it, so no unresolved reference or citation is produced. Printed slips kept literal: no printed word, symbol, hypothesis or number of the counted statement or of its proof was altered. Layout and class: the article's own class is \texttt{article} with packages including amsmath, amssymb, latexsym, mathtools, graphicx, epsfig, longtable, amsthm, xcolor, multirow, caption, subcaption, upgreek, makecell; the wrapper is the lane's pinned \texttt{amsart} editorial wrapper at 10pt on A4, which restates the theorem-like environments the retained text needs and adds the article's own macro definitions that the retained text needs (none), with extra wrapper packages (bm, bbm, dsfont, xspace, cancel, mathtools). The delivered numbering is the unit's own. Assets: no retained span contains a figure environment, a graphics-inclusion command, a PDF-page inclusion, a table or a TikZ picture; no image file is copied into this package, the package declares no asset dependency and no image file is redistributed with it. Licence: version arXiv:2505.22206v1 of this article is distributed under the Creative Commons Attribution 4.0 International licence, as recorded in the arXiv licence metadata for this exact version and shown on the abstract page https://arxiv.org/abs/2505.22206v1. A full rights notice is delivered at licenses/arxiv-2505.22206-CC-BY-4.0.txt. Characters: every retained excerpt is pure ASCII, so the delivered engine, XeLaTeX with its default Latin Modern text font, reports no missing character.
\begin{lemma}\label{lem:aux}
\label{lem:generalized_product_decomposition}
Let $U_1, \ldots, U_d$ be random variables with marginal distributions uniform on $[0,1]$. For any partition of the set of indices $\{1, \ldots, d\}$ into two disjoint sets $I$ and $J$ (i.e., $I \cup J = \{1, \ldots, d\}$ and $I \cap J = \emptyset$), the following identity holds for the expectation of a mixed product:
$$ \mathbb{E}\left[ \prod_{i \in I} U_i \prod_{i \in J} (1-U_i) \right] = \sum_{S \subseteq J} (-1)^{|S|} \mathbb{E}\left[ \prod_{l \in I \cup S} (1 - U_l) \right]. $$
This identity also applies to empirical averages for a sample $U_{1j}, \ldots, U_{dj}$ for $j=1, \ldots, n$:
$$ \frac{1}{n}\sum_{j=1}^n \left(\prod_{i \in I} U_{ij}\right) \left(\prod_{i \in J} (1-U_{ij})\right) = \sum_{S \subseteq J} (-1)^{|S|} \left( \frac{1}{n} \sum_{j=1}^n \prod_{l \in I \cup S} (1 - U_{lj}) \right). $$
\end{lemma}

\begin{theorem}{\label{th:est}}
    Let $\{X_{1j},\ldots,X_{d_j}\}_{j=1}^n$ be a $d$-dimensional random sample of size $n$ of a $d$-dimensional random vector $(X_1,\ldots,X_d)$ with associated $d$-copula $C$. Let $\alpha\in\mathbb{R}^n$ be such that $\alpha_i\in\{-1,1\}$ for all $i=1,\ldots,n$. Let $I\subseteq\{1,2,\ldots,n\}$ such that $\alpha_i=-1$ if $i\in I$ and $\alpha_i=1$ if $i\in J=\{1,2,\ldots,n\}\setminus I$% with $I,J\neq \emptyset$
    . Then
    %  \begin{align*}
    % \widehat{\rho}_d^\alpha(C)=&\frac{2^n(n+1)}{2^n-
    %(n+1)}\left(\frac{2^{|I|}-(|I|+1)}{2^{|I|}
    %(|I|+1)}\widehat{\rho}_{X_I}^- -\frac{2^{|I|+1}-(|I|+2)}{2^{|I|+1}
    %(|I|+2)}\sum_{j\in J}\widehat{\rho}_{X_{I\cup\{j\}}}^-\right.\\
    %     &\left.+\frac{2^{|I|+2}-(|I|+3)}{2^{|I|+2}
    %(|I|+3)}\sum_{j_1\in J}\sum_{\substack{j_2\in J\\ %j_2>j_1}}\widehat{\rho}_{X_{I\cup\{j_1,j_2\}}}^- -\cdots + 
    %(-1)^{|J|}\frac{2^n-(n+1)}{2^n(n+1)}\widehat{\rho}_X^-\right)
    % \end{align*}
\begin{align}\label{eq:estimator}
\widehat{\rho}_d^{\alpha} = \frac{(n+1)^d}{\frac{1}{n}\sum_{j=1}^nj^d-\left(\frac{n+1}{2}\right)^d} \sum_{k=0}^{|J|} (-1)^k \frac{2^{|I|+k}-(|I|+k+1)}{2^{|I|+k}(|I|+k+1)} \sum_{\substack{S \subseteq J \\ |S|=k}} \widehat{\rho}_{X_{I\cup S}}^-
\end{align}
\end{theorem}

\begin{proof}Let $\mathbf{X}_1, \ldots, \mathbf{X}_n$ be a random sample from a $d$-variate distribution, and let $R_{ij}^{\alpha_i}$ the (directional) ranks. The pseudo-observations are defined as $U_{ij} = R_{ij}/(n+1)$ for $i=1,\ldots,d$ and $j=1,\ldots,n$. For a subset of indices $K \subseteq \{1,\ldots,d\}$ with $|K|$ components, the empirical upper-tail dependence coefficient $\widehat{\rho}_K^-$ is defined as:
$$\widehat{\rho}_K^- = \frac{2^{|K|}(|K|+1)}{2^{|K|}-(|K|+1)} \left( \frac{1}{n} \sum_{j=1}^n \prod_{i \in K} (1 - U_{ij}) - \left(\frac{1}{2}\right)^{|K|} \right).$$
Now, let $N_{\text{rank}}$ denote the numerator of $\widehat{\rho}_d^{\alpha}$ and $D_{\text{rank}}$ denote its denominator, i.e.,
$$N_{\text{rank}} = \frac{1}{n}\sum_{j=1}^n\prod_{i=1}^dR_{ij}^{\alpha_i}-\left(\frac{n+1}{2}\right)^d$$
$$D_{\text{rank}} = \frac{1}{n}\sum_{j=1}^nj^d-\left(\frac{n+1}{2}\right)^d$$
By definition, $R_{ij}^{\alpha_i} = (n+1)U_{ij}$ if $\alpha_i = -1$ (i.e., $i \in I$), and $R_{ij}^{\alpha_i} = (n+1)(1-U_{ij})$ if $\alpha_i = 1$ (i.e., $i \in J$). Thus, the product $\prod_{i=1}^d R_{ij}^{\alpha_i}$ becomes:
\begin{align*} \prod_{i=1}^d R_{ij}^{\alpha_i} &= \left(\prod_{i \in I} (n+1)U_{ij}\right) \left(\prod_{i \in J} (n+1)(1-U_{ij})\right) \\ &= (n+1)^{|I|} (n+1)^{|J|} \left(\prod_{i \in I} U_{ij}\right) \left(\prod_{i \in J} (1-U_{ij})\right) \\ &= (n+1)^d \left(\prod_{i \in I} U_{ij}\right) \left(\prod_{i \in J} (1-U_{ij})\right).
\end{align*}
Substituting this expression into the expression for $N_{\text{rank}}$:
\begin{align*} N_{\text{rank}} &= \frac{1}{n}\sum_{j=1}^n (n+1)^d \left(\prod_{i \in I} U_{ij}\right) \left(\prod_{i \in J} (1-U_{ij})\right) - \left(\frac{n+1}{2}\right)^d \\ &= (n+1)^d \left[ \frac{1}{n}\sum_{j=1}^n \left(\prod_{i \in I} U_{ij}\right) \left(\prod_{i \in J} (1-U_{ij})\right) - \left(\frac{1}{2}\right)^d \right]. \end{align*}
Let $K_{\alpha}^{\text{emp}}$ denote the term in the square brackets:
$$K_{\alpha}^{\text{emp}} = \frac{1}{n}\sum_{j=1}^n \left(\prod_{i \in I} U_{ij}\right) \left(\prod_{i \in J} (1-U_{ij})\right) - \left(\frac{1}{2}\right)^d;$$
so, $N_{\text{rank}} = (n+1)^d K_{\alpha}^{\text{emp}}$. Lemma \ref{lem:aux} states that $K_{\alpha}^{\text{emp}}$ can be expressed as:
$$K_{\alpha}^{\text{emp}} = \sum_{S \subseteq J} (-1)^{|S|} \left( \frac{1}{n} \sum_{j=1}^n \prod_{l \in I \cup S} (1 - U_{lj}) - \left(\frac{1}{2}\right)^{|I|+|S|} \right).$$
This identity allows us to decompose $K_{\alpha}^{\text{emp}}$ into a sum of terms related to products of $(1-U_{lj})$.
From the definition of $\widehat{\rho}_K^-$, we can isolate the parenthesized term:
$$ \frac{1}{n} \sum_{j=1}^n \prod_{i \in K} (1 - U_{ij}) - \left(\frac{1}{2}\right)^{|K|} = \frac{2^{|K|}-(|K|+1)}{2^{|K|}(|K|+1)} \widehat{\rho}_K^-. $$
Substituting this, and by setting $K = I \cup S$, we have:
$$K_{\alpha}^{\text{emp}} = \sum_{S \subseteq J} (-1)^{|S|} \left( \frac{2^{|I|+|S|}-(|I|+|S|+1)}{2^{|I|+|S|}(|I|+|S|+1)} \widehat{\rho}_{X_{I\cup S}}^- \right).$$
This sum can be regrouped by the size of $S$, $k=|S|$:
$$K_{\alpha}^{\text{emp}} = \sum_{k=0}^{|J|} (-1)^k \frac{2^{|I|+k}-(|I|+k+1)}{2^{|I|+k}(|I|+k+1)} \sum_{\substack{S \subseteq J \\ |S|=k}} \widehat{\rho}_{X_{I\cup S}}^-.$$
Finally, recall that $\widehat{\rho}_d^{\alpha} = \frac{N_{\text{rank}}}{D_{\text{rank}}}$ and substitute the expression for $N_{\text{rank}} = (n+1)^d K_{\alpha}^{\text{emp}}$, whence Equation \eqref{eq:estimator} follows, completing the proof.
%Since the estimators $\widehat{\rho}_d^-$ are unbiased for every %dimension (see Subsection \ref{sect:prop}), then
%\begin{align*}
%    &\mathbb{E}\left[\frac{2^n(n+1)}{2^n-(n+1)}\left(\frac{2^{|I|}-
%(|I|+1)}{2^{|I|}(|I|+1)}\widehat{\rho}_{X_I}^- -\frac{2^{|I|+1}-
%(|I|+2)}{2^{|I|+1}(|I|+2)}\sum_{j\in J}\widehat{\rho}_{X_{I\cup\
%{j\}}}^-\right.\right.\\
%    &\left.\left.+\frac{2^{|I|+2}-(|I|+3)}{2^{|I|+2}
%(|I|+3)}\sum_{j_1\in J}\sum_{\substack{j_2\in J\\ %j_2>j_1}}\widehat{\rho}_{X_{I\cup\{j_1,j_2\}}}^- -\cdots + 
%(-1)^{|J|}\frac{2^n-(n+1)}{2^n(n+1)}\widehat{\rho}_X^-\right)\right]\\
%    =&\frac{2^n(n+1)}{2^n-(n+1)}\left(\frac{2^{|I|}-(|I|+1)}{2^{|I|}
%(|I|+1)}\mathbb{E}[\widehat{\rho}_{X_I}^-] -\frac{2^{|I|+1}-(|I|+2)}
%{2^{|I|+1}(|I|+2)}\sum_{j\in J}\mathbb{E}[\widehat{\rho}_{X_{I\cup\
%{j\}}}^-]\right.\\
%    &\left.+\frac{2^{|I|+2}-(|I|+3)}{2^{|I|+2}(|I|+3)}\sum_{j_1\in %J}\sum_{\substack{j_2\in J\\ j_2>j_1}}\mathbb{E}
%[\widehat{\rho}_{X_{I\cup\{j_1,j_2\}}}^-] -\cdots + 
%(-1)^{|J|}\frac{2^n-(n+1)}{2^n(n+1)}\mathbb{E}
%[\widehat{\rho}_X^-]\right)\\
%    =&\widehat{\rho}_d^\alpha.
%\end{align*}
\end{proof}

\end{document}
